\section*{Sets and Indices}

\begin{itemize}
    \item $R$: set of candidate restaurants, indexed by $r$.
    \item In this instance,
    \[
    R=\{\text{Restaurant\_A},\text{Restaurant\_B},\text{Restaurant\_C},\text{Restaurant\_D}\}.
    \]
\end{itemize}

\section*{Parameters}

\begin{itemize}
    \item $c_r$: investment cost of restaurant $r$.
    \item $v_r$: annual revenue of restaurant $r$.
    \item $B$ (code: \texttt{budget}): total investment budget.
\end{itemize}

Using the provided data:
\[
B = 6{,}000{,}000.
\]

For each restaurant:
\[
\begin{array}{c|cc}
r & v_r & c_r \\
\hline
\text{Restaurant\_A} & 15{,}000 & 1{,}600{,}000 \\
\text{Restaurant\_B} & 40{,}000 & 2{,}500{,}000 \\
\text{Restaurant\_C} & 30{,}000 & 1{,}800{,}000 \\
\text{Restaurant\_D} & 50{,}000 & 3{,}000{,}000 \\
\end{array}
\]

\section*{Decision Variables}

\begin{itemize}
    \item $x_r$ (code: \texttt{x[r['name']]}): binary variable equal to $1$ if restaurant $r$ is purchased, and $0$ otherwise.
\end{itemize}

Thus,
\[
x_r \in \{0,1\} \qquad \forall r \in R.
\]

\section*{Objective}

Maximize total annual revenue:
\[
\max \sum_{r \in R} v_r x_r.
\]

For this instance:
\[
\max\;
15000\,x_{\text{Restaurant\_A}}
+40000\,x_{\text{Restaurant\_B}}
+30000\,x_{\text{Restaurant\_C}}
+50000\,x_{\text{Restaurant\_D}}.
\]

\section*{Constraints}

\begin{align}
\sum_{r \in R} c_r x_r &\le B
&& \text{(budget)} \\
x_{\text{Restaurant\_D}} + x_{\text{Restaurant\_A}} &\le 1
&& \text{(D--A exclusion)} \\
x_r &\in \{0,1\}
&& \forall r \in R.
\end{align}

For this instance, the budget constraint is:
\[
1600000\,x_{\text{Restaurant\_A}}
+2500000\,x_{\text{Restaurant\_B}}
+1800000\,x_{\text{Restaurant\_C}}
+3000000\,x_{\text{Restaurant\_D}}
\le 6000000.
\]

\section*{Notes on Implementation}

\begin{itemize}
    \item The model is implemented as a binary linear optimization problem in PuLP-compatible syntax and solved with \texttt{GUROBI\_CMD(msg=0)}.
    \item Restaurant data are loaded from \texttt{table\_1.csv} into records with fields \texttt{name}, \texttt{revenue}, and \texttt{cost}; the budget is loaded from \texttt{general\_parameters.csv} under \texttt{investment\_budget}.
    \item The exclusion constraint between Restaurant D and Restaurant A is added only if both keys \texttt{'Restaurant\_D'} and \texttt{'Restaurant\_A'} are present in the loaded data. For this instance, both are present, so the constraint
    \[
    x_{\text{Restaurant\_D}} + x_{\text{Restaurant\_A}} \le 1
    \]
    is part of the solved model.
    \item The code does not use the CSV parameter \texttt{restaurant\_d\_a\_constraint} numerically; instead, the D--A exclusion is hard-coded through the conditional check described above.
\end{itemize}
